% Abstract (150 words)

Congressional redistricting in the United States operates under a constitutional constraint: districts cannot cross state boundaries. This paper quantifies the geometric consequences of this federalist structure through a computational experiment. We apply graph partitioning algorithms to redistrict all 435 House seats nationally, treating the United States as a single graph of 84,414 census tracts with no regard for state boundaries. Comparing national optimization to state-based redistricting reveals a counterintuitive finding: national optimization produces districts 42.5\% \textit{less compact} than state-based redistricting (Polsby-Popper scores of 0.265 vs 0.461). This result challenges conventional assumptions that geographic constraints harm optimization objectives. An ablation study varying boundary crossing penalties confirms that state boundaries facilitate rather than impede geometric compactness. While 41.4\% of nationally optimized districts cross state lines, suggesting fractional representation as a theoretical alternative, no configuration improves upon state-based compactness. These findings demonstrate that federalist redistricting structures provide beneficial organization for optimization algorithms rather than imposing geometric costs.

\paragraph{Comparability Note} The baseline PP of 0.461 in this paper reflects the \textit{national redistricting configuration} (no state boundaries, 435 contiguous districts across the continental US evaluated against the state-based control). This figure differs from the per-state B.2 result (mean PP = 0.361) because the national partition produces more geometrically compact districts when states are used as the unit of analysis --- states like Wyoming and Montana, which have poor compactness when forced into a single district, benefit from being merged with neighboring areas in the national configuration. The two configurations are not directly comparable: B.2 reports per-state means across all 50 states individually, while this paper's 0.461 is the aggregate Polsby-Popper for the state-based control run (all 50 states partitioned simultaneously as independent subproblems in a single experiment).
